% !TeX program = xelatex
% =============================================================
% 模板 10：IEEE 期刊投稿（IEEEtran journal，双栏 10pt）
% 编译： scripts\build.ps1 en-journal-ieee
% 版式： 双栏 10pt，IEEE 官方 IEEEtran.cls
% 用途： IEEE Transactions / Journals 初稿；会议投稿用 en-ieee
% =============================================================
\documentclass[journal,10pt]{IEEEtran}

\usepackage{cite}
\usepackage{amsmath}
\usepackage{fontspec}
% IEEEtran 在 XeLaTeX 下不会自动加载 Times，这里显式换成度量兼容克隆
\setmainfont{TeX Gyre Termes}
\usepackage{unicode-math}
\setmathfont{TeX Gyre Termes Math}
\usepackage{amsfonts}
\usepackage{graphicx}
\usepackage{booktabs}
\usepackage{multirow}
\usepackage{textcomp}
\usepackage{xcolor}
\usepackage[hidelinks]{hyperref}

\begin{document}

% ------------------------- 标题与作者（改这里） -------------------------
\title{Placeholder Title: A Sample IEEE Transactions Manuscript for Template Testing}

\author{San Zhang,~\IEEEmembership{Student Member,~IEEE,}
  Si Li,~\IEEEmembership{Member,~IEEE,}
  and Wu Wang,~\IEEEmembership{Senior Member,~IEEE}%
\thanks{Manuscript received 1 January 2026; revised 1 March 2026; accepted 1 April 2026.
  Date of publication 1 May 2026; date of current version 1 June 2026. This work was
  supported by the National Natural Science Foundation of China under Grant 00000000.
  (Corresponding author: Si Li.)}%
\thanks{San Zhang and Si Li are with the School of Electrical Engineering, XX University,
  Beijing 100000, China (e-mail: zhangsan@example.edu.cn; lisi@example.edu.cn).}%
\thanks{Wu Wang is with the School of Automation, YY University, Shanghai 200000, China
  (e-mail: wangwu@example.edu.cn).}}

\markboth{IEEE Transactions on Placeholder Studies, Vol. 1, No. 1, January 2026}%
{Zhang \MakeLowercase{\textit{et al.}}: Placeholder IEEE Transactions Manuscript}

\maketitle

% ------------------------- 摘要与关键词 -------------------------
\begin{abstract}
This manuscript is a placeholder for an IEEE Transactions submission. It demonstrates the
title block, the abstract, a numbered equation, a table, and a bibliography entry so that
the layout can be checked before the real content is written. Replace every sentence here
with your own 150--250 word abstract, which should state the problem, the approach, the
main quantitative results, and the conclusion, without citations or undefined abbreviations.
\end{abstract}

\begin{IEEEkeywords}
\raggedright
placeholder, template, IEEEtran, journal manuscript, two-column layout, formatting check
\end{IEEEkeywords}

% ------------------------- 正文 -------------------------
\section{Introduction}
\label{sec:intro}
This is the introduction placeholder. An IEEE Transactions introduction
usually states the problem, critically reviews related work~\cite{lecun2015deep,he2016resnet},
identifies the research gap, and ends with the contributions of this paper.

\section{Method}
\label{sec:method}
Let $x_i$ denote the $i$-th input feature and $w_i$ its weight. The weighted sum is
\begin{equation}
  y=\sum_{i=1}^{n} w_i x_i + b,
  \label{eq:linear}
\end{equation}
where $b$ is the bias term. Equation~\eqref{eq:linear} is the equation placeholder.

\section{Experimental Setup}
\label{sec:setup}
The parameter settings are listed in Table~\ref{tab:params}.

\begin{table}[h]
  \renewcommand{\arraystretch}{1.2}
  \caption{Placeholder Experimental Parameters}
  \label{tab:params}
  \centering
  \begin{tabular}{lcc}
    \toprule
    Parameter & Symbol & Value \\
    \midrule
    Learning rate & $\eta$    & 0.01 \\
    Epochs        & $E$       & 100  \\
    Batch size    & $B$       & 32   \\
    Weight decay  & $\lambda$ & 0.10 \\
    \bottomrule
  \end{tabular}
\end{table}

\section{Results and Discussion}
\label{sec:results}
With the settings in Table~\ref{tab:params}, the placeholder model reaches 92.4\% accuracy
on the validation split. Report mean and standard deviation over repeated runs, state the
statistical test used, and discuss why the improvement holds rather than only reporting it.

\section{Conclusion}
\label{sec:conclusion}
This template compiles end-to-end with XeLaTeX and BibTeX. Replace the placeholder text,
figures, and references with your own content before submission. Check the target
Transactions' Information for Authors for page limits, reference style, and supplementary
material requirements.

% ------------------------- 附录（可选，不需要就删掉整段） -------------------------
\appendices
\section{Derivation of the Weighted Sum}
\label{app:derivation}
Expand (1) term by term to obtain the full derivation used in Section~\ref{sec:method}.

% ------------------------- 参考文献 -------------------------
\bibliographystyle{IEEEtran}
\bibliography{refs}

% ------------------------- 作者简介（可选，投稿时按需填） -------------------------
% \begin{IEEEbiography}[{\includegraphics[width=1in,height=1.25in,clip,keepaspectratio]{photo1}}]{San Zhang}
% Bio text here.
% \end{IEEEbiography}

\end{document}
